\begin{solapartat}
\punts{0,3} Podem comprovar que la primera freqüència que apareix, l'harmònic fonamental, correspon a $f_1 = 440\ \mathrm{Hz}$.

\punts{0,65} Segons les condicions de frontera, les ones estacionàries són per a un piano i un oboè:

\figura[0.9]{sol_fig1.png}

Podem comprovar, doncs, que per a un piano:
\[ L = n\frac{\lambda}{2} \Rightarrow f = n\frac{v}{2L} = nf_1, \qquad n = 1, 2, 3, \ldots \]

I per a un oboè:
\[ L = n\frac{\lambda}{4} \Rightarrow f = n\frac{v}{4L} = nf_1, \qquad n = 1, 3, 5, \ldots \]

Si es tractés d'un piano, el segon harmònic seria: $f_2 = 2\cdot 440 = 880\ \mathrm{Hz}$, mentre que en un oboè no apareix el segon harmònic. El primer harmònic que apareix en un oboè és el tercer: $f_3 = 3\cdot 440 = 1320\ \mathrm{Hz}$.

A partir de l'espectre de Fourier veiem que la segona freqüència que apareix és 1320 Hz, per tant, l'espectre no pot correspondre a un piano, així la resposta és \textbf{oboè.}

\punts{0,3} La freqüència fonamental per a un oboè és:
\[ f_1 = \frac{v}{4L} \Rightarrow L = \frac{v}{4f_1} = \frac{340}{4\cdot 440} = 0,193\ \mathrm{m} \]
\end{solapartat}

\begin{solapartat}
\punts{0,65} Intensitat sonora que ens arriba de cada instruments $I$ és,
\[ \beta = 10\log\frac{5I}{I_0} \Rightarrow \frac{\beta}{10} = \log\frac{5I}{I_0} \Rightarrow I = \frac{I_0 10^{\beta/10}}{5} = 2,00\times 10^{-5}\ \mathrm{W\,m^{-2}} \]

\punts{0,6}
\[ I = \frac{\text{Potència}}{A} \Rightarrow \text{Potència} = I\,2\pi r^2 = 2,83\times 10^{-4}\ \mathrm{W} \]
\end{solapartat}
